%%%%%%%%%%%%%%%%%%%%%%%%%%%%%%%%%%%%%%%%%%%%%%%%%%%%%%%%%%%%%%%%%%%%%%%%%%%%%%%%

\pagestyle{empty}
\begin{abstract}
  En el desarrollo moderno de videojuegos, la navegación autónoma constituye la base de la inteligencia artificial que dota de credibilidad a los personajes no jugables.
  Las soluciones más extendidas, basadas en las llamadas mallas de navegación, están diseñadas para agentes terrestres que se desplazan sobre superficies.
  Este enfoque deja sin resolver un caso de uso de gran importancia, el de los agentes voladores como drones, naves o criaturas que necesitan moverse libremente a través de la totalidad del espacio libre.

  Para cubrir este vacío, este proyecto presenta un sistema de navegación volumétrica completo, eficiente y de código abierto, integrado en el motor Unity.
  El resultado no es únicamente un algoritmo de búsqueda de rutas, sino un paquete de herramientas modular y listo para producción.
  Su integración nativa en dicho motor, con Gizmos de depuración visual y paneles de configuración, hace que el sistema resulte accesible tanto para perfiles técnicos como para diseñadores de niveles y artistas.
  
  \vspace*{25pt}

  \begin{segundoresumo}
  In modern video game development, autonomous navigation forms the foundation of the artificial intelligence that gives credibility to non-playable characters.
  The most widespread solutions, based on so-called navigation meshes, are designed for ground-based agents moving across surfaces.
  This approach leaves a highly important use case unresolved: flying agents such as drones, ships, or creatures that need to move freely throughout the entirety of open space.

  To fill this gap, this project presents a complete, efficient, and open-source volumetric navigation system integrated into the Unity engine.
  The result is not merely a pathfinding algorithm, but a modular, production-ready toolset.
  Its native integration into the engine, complete with visual debugging Gizmos and configuration panels, makes the system accessible to technical profiles as well as level designers and artists.

  \end{segundoresumo}
\vspace*{25pt}
\begin{multicols}{2}
\begin{description}
\item [\palabraschaveprincipal:] \mbox{} \\[-20pt]
  \begin{itemize}
    \item Navegación volumétrica
    \item A*
    \item Pathfinding
    \item Path planning
    \item Heurísticas
    \item Vóxel
    \item Octree
    \item Entorno inmersivo
    \item Unity
  \end{itemize}
\end{description}
\begin{description}
\item [\palabraschavesecundaria:] \mbox{} \\[-20pt]
  \begin{itemize}
    \item Volumetric navigation
    \item A*
    \item Pathfinding
    \item Path planning
    \item Heuristics
    \item Voxel
    \item Octree
    \item Immersive environment
    \item Unity
  \end{itemize}
\end{description}
\end{multicols}

\end{abstract}
\pagestyle{fancy}

%%%%%%%%%%%%%%%%%%%%%%%%%%%%%%%%%%%%%%%%%%%%%%%%%%%%%%%%%%%%%%%%%%%%%%%%%%%%%%%%
